\documentclass[11pt]{article}
\usepackage[margin=1in]{geometry}
\usepackage{amsmath,amssymb,amsthm}
\usepackage{graphicx}
\usepackage{booktabs}
\usepackage{listings}
\usepackage{xcolor}
\usepackage[numbers]{natbib}
\usepackage{hyperref}

\lstset{basicstyle=\ttfamily\footnotesize,breaklines=true,columns=fullflexible,
  frame=single,extendedchars=true,
  literate={ℝ}{{$\mathbb{R}$}}1 {∀}{{$\forall$}}1 {→}{{$\to$}}1 {≤}{{$\le$}}1
    {↔}{{$\leftrightarrow$}}1 {∧}{{$\wedge$}}1 {·}{{$\cdot$}}1 {−}{{$-$}}1
    {–}{{--}}1 {≥}{{$\ge$}}1 {¬}{{$\neg$}}1 {≠}{{$\ne$}}1 {↦}{{$\mapsto$}}1
    {∈}{{$\in$}}1 {×}{{$\times$}}1 {⟨}{{$\langle$}}1 {⟩}{{$\rangle$}}1
    {`}{{\textasciigrave}}1}

\newtheorem{theorem}{Theorem}
\newtheorem{lemma}{Lemma}

\title{The Levins extinction threshold under habitat loss:\\ a Lean proof of the sign of the growth term}
\author{Justin Hart\\Viridis LLC}
\date{28 November 2026}

\begin{document}
\maketitle

\begin{abstract}
In the Levins metapopulation model with habitat loss, the fraction $p$ of patches occupied
changes at rate $c\,p\,(h-p)-e\,p$, where $c$ is the colonisation rate, $e$ the extinction rate
and $h$ the fraction of patches still suitable. The classical conclusion
\citep{Levins1969,Lande1987,NeeMay1992} is that a population persists only while $h$ exceeds the
threshold $e/c$. Equivalently, the destroyed fraction $1-h$ must stay below $1-e/c$, the
occupancy the species had in intact habitat. This paper continues project P7 of Viridis Research.
We give a Lean~4 proof, over Mathlib, of the sealed statement \texttt{P7.HabitatThreshold}. The
statement has five parts: a factorisation of the growth term through the equilibrium
$p^*=h-e/c$; a negative sign at or below the threshold; the location and attracting sign pattern
of $p^*$ above it; the Levins rule; and strict monotonicity of the threshold in $c$ and $e$. What
the proof settles is the sign of the growth term in this model. It says nothing about solutions
of the differential equation, nothing about whether any species follows the model, and nothing
about any value of $c$, $e$ or $h$.
\end{abstract}

\section{Introduction}
How much habitat does a population need? The Levins metapopulation model \citep{Levins1969},
extended to habitat destruction \citep{Lande1987,NeeMay1992}, answers with a threshold: a species
occupying a fraction $p$ of patches declines whenever the suitable fraction $h$ of patches falls
to $e/c$ or below. Here $e$ is the per-patch extinction rate and $c$ the colonisation rate. In
intact habitat ($h=1$) the equilibrium occupancy is $1-e/c$. The threshold can therefore be
stated as the \emph{Levins rule}: the species persists exactly while the destroyed fraction $1-h$
is less than the occupancy it had in intact habitat.

The result is classical and the algebra is elementary. A habitat-loss kernel that relies on it
should still state precisely which facts it uses and under which hypotheses, and those facts
should be checked by machine. The previous session of project P7 (P7-T001) proposed the statement
\texttt{P7.HabitatThreshold}, which has since been sealed (decision D-027). Tonight's task, P7-T002,
is to prove it. We give a Lean proof that compiles under the pinned toolchain and depends only on
the standard axioms \texttt{propext}, \texttt{Classical.choice} and \texttt{Quot.sound}.

\section{Background and related work}
\paragraph{The model.} \citet{Levins1969} introduced the patch-occupancy model
$dp/dt=c\,p\,(1-p)-e\,p$ for a population spread over many identical patches, each either occupied
or empty. Colonisation of empty patches is proportional to the occupied fraction, and extinction
is at a constant per-patch rate. \citet{Lande1987} studied extinction thresholds in demographic
models of territorial populations, in which only a fraction $h$ of sites is suitable.
\citet{NeeMay1992} added habitat destruction to the Levins model in the form
\[ \frac{dp}{dt} = c\,p\,(h-p) - e\,p = c\,p\,(p^*-p),\qquad p^* = h - \frac{e}{c}, \]
and observed that the equilibrium occupancy falls one-for-one with the destroyed fraction. The
species is lost once $1-h$ reaches $1-e/c$.

\paragraph{The companion statement.} P7 also includes \texttt{P7.SpeciesArea}, which bounds the
fraction of species kept under a power-law species--area relationship \citep{Arrhenius1921}. It
was proposed together with this statement, and its proof is task P7-T003. It is not treated here.

\paragraph{Connectivity.} Corridor networks, the subject of project P6 of this programme, enter
the Levins model only through the colonisation rate $c$. Better connectivity, if it raises $c$,
lowers the threshold $e/c$ (part~5 below), and that is the only route by which a network appears
in this statement. The paper makes no claim about how any corridor network changes $c$, by how
much, or whether it does so at all. That is an empirical question about dispersal, and nothing
proved here bears on it.

\section{Formal development}
The sealed module defines, over $\mathbb{R}$,
\[ \mathrm{growth}(c,e,h,p)=c\,p\,(h-p)-e\,p,\qquad p^*(c,e,h)=h-e/c. \]

\begin{theorem}[\texttt{P7.HabitatThreshold}]
For all real $c,e,h$ with $c>0$ and $e>0$:
\begin{enumerate}
\item $\mathrm{growth}(p)=c\,p\,(p^*-p)$ for every real $p$;
\item if $h\le e/c$, then $\mathrm{growth}(p)<0$ for every $p>0$;
\item if $e/c<h$, then $0<p^*<h$, $\mathrm{growth}(p^*)=0$, $\mathrm{growth}(p)>0$ for
$0<p<p^*$, and $\mathrm{growth}(p)<0$ for $p>p^*$;
\item $e/c<h \iff 1-h<1-e/c$;
\item if $c<c'$ then $e/c'<e/c$, and if $e<e'$ then $e/c<e'/c$.
\end{enumerate}
\end{theorem}

\begin{proof}[Proof sketch]
(1) Since $c\neq0$, $c\,p\,(h-e/c-p)=c\,p\,(h-p)-e\,p$. (2) By (1), $\mathrm{growth}(p)$ is the
product of $c\,p>0$ and $p^*-p$. The second factor is negative because $p^*=h-e/c\le0<p$.
(3) $p^*>0$ is the hypothesis rearranged, and $p^*<h$ because $e/c>0$, which is where $e>0$ is
used. At $p=p^*$ the second factor of (1) vanishes. For $0<p<p^*$ both factors are positive. For
$p>p^*>0$ the first factor is positive and the second negative. (4) Both sides are the same linear
inequality. (5) Division of a positive number by a larger positive number gives a smaller
quotient, and division by a fixed positive number is strictly increasing.
\end{proof}

The Lean proof follows this sketch. One helper lemma proves the factorisation with
\texttt{field\_simp} and \texttt{ring}. Parts~2 and~3 rewrite with it and close the sign goals with
\texttt{mul\_pos} and \texttt{mul\_neg\_of\_pos\_of\_neg}, using \texttt{linarith} for the
side conditions. Part~5 uses the Mathlib lemmas \texttt{div\_lt\_div\_of\_pos\_left} and
\texttt{div\_lt\_div\_of\_pos\_right}. The complete candidate file is:

\begin{lstlisting}
import Mathlib
import ViridisStatements.P7.HabitatThreshold

namespace Candidate

open ViridisStatements.P7.HabitatThreshold

/-- Part 1: the growth term factorises through the equilibrium `p* = h - e/c`. -/
theorem growth_factor (c e h : ℝ) (hc : 0 < c) (p : ℝ) :
    growth c e h p = c * p * (pstar c e h - p) := by
  unfold growth pstar
  field_simp
  ring

/-- Part 2: at or below the threshold the growth term is negative for every `p > 0`. -/
theorem below (c e h : ℝ) (hc : 0 < c) (hle : h ≤ e / c) (p : ℝ) (hp : 0 < p) :
    growth c e h p < 0 := by
  rw [growth_factor c e h hc p]
  have hcp : 0 < c * p := mul_pos hc hp
  have hneg : pstar c e h - p < 0 := by unfold pstar; linarith
  exact mul_neg_of_pos_of_neg hcp hneg

/-- Part 3: above the threshold, `0 < p* < h`, `p*` is an equilibrium, and the growth term
points towards `p*` from both sides. -/
theorem above (c e h : ℝ) (hc : 0 < c) (he : 0 < e) (hlt : e / c < h) :
    0 < pstar c e h ∧ pstar c e h < h ∧ growth c e h (pstar c e h) = 0 ∧
      (∀ p : ℝ, 0 < p → p < pstar c e h → 0 < growth c e h p) ∧
      (∀ p : ℝ, pstar c e h < p → growth c e h p < 0) := by
  have hec : 0 < e / c := div_pos he hc
  have hps : 0 < pstar c e h := by unfold pstar; linarith
  refine ⟨hps, ?_, ?_, ?_, ?_⟩
  · unfold pstar; linarith
  · rw [growth_factor c e h hc]; ring
  · intro p hp hpp
    rw [growth_factor c e h hc p]
    exact mul_pos (mul_pos hc hp) (by linarith)
  · intro p hpp
    rw [growth_factor c e h hc p]
    exact mul_neg_of_pos_of_neg (mul_pos hc (by linarith)) (by linarith)

theorem main : ViridisStatements.P7.HabitatThreshold.Statement := by
  refine ⟨?_, ?_, ?_, ?_, ?_⟩
  · intro c e h hc _ p
    exact growth_factor c e h hc p
  · intro c e h hc _ hle p hp
    exact below c e h hc hle p hp
  · intro c e h hc he hlt
    exact above c e h hc he hlt
  · intro c e h _ _
    constructor <;> intro H <;> linarith
  · intro c e hc he
    refine ⟨?_, ?_⟩
    · intro c' hc'
      exact div_lt_div_of_pos_left he hc hc'
    · intro e' he'
      exact div_lt_div_of_pos_right he' hc

end Candidate
\end{lstlisting}

\section{Results}
The candidate file \texttt{proofs/P7.HabitatThreshold.lean} passes all five stages of the
programme's local checker \texttt{viridis-lean-check} (denylist, compilation, statement match,
axioms, kernel replay). The only axioms reported are \texttt{propext}, \texttt{Classical.choice}
and \texttt{Quot.sound}. It passed on the first attempt, and no automated prover was needed. The
certification page at the front of this paper records the verifier's own result.

We also observe how the hypotheses are used. The assumption $c>0$ is used in parts 1--3 and 5.
The assumption $e>0$ is used only for $p^*<h$ in part~3 and for strictness in part~5. Part~4 is
pure linear arithmetic and needs neither. The sealed module's negative controls show that the
hypotheses $p>0$ in part~2 and $e>0$ in part~5 cannot be dropped.

\section{What this settles and what it does not}
\paragraph{Settled.} In the model $dp/dt=c\,p\,(h-p)-e\,p$ with $c,e>0$, the sign of the growth term
is as the classical analysis says. At or below $h=e/c$ it pushes every positive occupancy down.
Above it there is exactly one positive equilibrium $p^*=h-e/c$, and the growth term points
towards $p^*$ from both sides. The threshold condition is the Levins rule, and the threshold
moves in the stated directions with $c$ and $e$.

\paragraph{Not settled.} The statement says nothing about solutions of the differential
equation: not their existence, not convergence to $p^*$, not how fast occupancy declines below
the threshold. Each of those would need a separate statement about trajectories. It does not
claim that any real species follows the Levins model. The model assumes identical patches,
global (mean-field) dispersal, constant rates, an infinite number of patches without demographic
stochasticity, and habitat loss that removes patches at random. It gives no value of $c$, $e$ or
$h$ for any species or landscape. Those are field inputs. It also does not address extinction
debt, the time lag before a population above or below the threshold settles. The parameters $p$
and $h$ are not even constrained to $[0,1]$ in the statement, since the algebra does not need it.

\section{Limitations}
The result is a classical one, restated and checked; it is not new mathematics. Its use for
planning depends entirely on the modelling premises listed in \texttt{run.json} (mean-field Levins
dynamics with constant rates, and random habitat loss), which no proof can establish.

On the programme's frontier duty for premise \texttt{M-Nucleation} (classical nucleation for
planted tree clusters), we again did not propose a test plan. A deciding test needs the size of
each planted cluster at two or more times, across initial sizes that bracket the predicted
critical size. The public data we could identify from the Costa Rica applied-nucleation
experiment \citep{Zahawi2013} and from a New Jersey landfill experiment \citep{Corbin2016} either
records recruitment rather than the size of each island over time, or uses a single planted plot
size. Neither could come out contradicted on the premise's main content: a single sign change in
cluster growth at a critical size. Details are in the notebook.

\section{Reproducibility}
The proof is \texttt{proofs/P7.HabitatThreshold.lean}. It targets the sealed module
\texttt{ViridisStatements.P7.HabitatThreshold} (sha256
\texttt{378845bc63602cb8290b25faa64b3dbff5b76cabcc6e2f33706a1551f6b431e8}) and compiles under
\texttt{leanprover/lean4:v4.28.0} with Mathlib commit
\texttt{8f9d9cff6bd728b17a24e163c9402775d9e6a365}. Check it with
\texttt{viridis-lean-check proofs/P7.HabitatThreshold.lean P7.HabitatThreshold}. No data or
analysis code was used.

\section{Changelog}
\begin{itemize}
\item 2026-11-27 (P7-T001): first version. Proposed \texttt{P7.HabitatThreshold} and
\texttt{P7.SpeciesArea}, with sanity checks, negative controls and rationales. No proofs.
\item 2026-11-28 (P7-T002): \texttt{P7.HabitatThreshold} sealed (D-027). Gave a Lean proof of it.
Rewrote the paper around the proof and the scope of what it settles. Added the paragraph on
corridor networks and $c$, and a premise review. The test plan for \texttt{M-Nucleation} is still
not proposed (reasons updated). Material on \texttt{P7.SpeciesArea} is deferred to P7-T003.
\end{itemize}

\bibliographystyle{plainnat}
\bibliography{refs}
\end{document}
